\documentclass[11pt]{article}
\usepackage{amsmath,amssymb,amsthm}
\usepackage[hidelinks]{hyperref}
\newtheorem{theorem}{Theorem}
\newcommand{\ket}[1]{\left|#1\right\rangle}

\title{Hello, World}
\author{Adam Holmes}
\date{\today}

\begin{document}
\maketitle

\begin{abstract}
A first post, written in \emph{folio} and typeset straight onto the page.
\end{abstract}

This blog is written the way I like to write papers: plain \LaTeX{} on disk, rendered in place
while I type. Inline math like $\ket{\psi} = \alpha\ket{0} + \beta\ket{1}$ sits in the prose,
and displayed equations are numbered for free:
\begin{equation}
  |\alpha|^2 + |\beta|^2 = 1 .
  \label{eq:norm}
\end{equation}
Equation~\eqref{eq:norm} is just normalisation.\footnote{Footnotes become endnotes on the web.}

\section{What to expect}
Short notes on quantum error correction, compilers and architecture --- things too small
for a paper and too long for a tweet~\cite{shor1995}.

\begin{theorem}[No-cloning]
\label{thm:nc}
There is no unitary $U$ with $U\ket{\psi}\ket{0} = \ket{\psi}\ket{\psi}$ for every $\ket{\psi}$.
\end{theorem}
\begin{proof}
Inner products are preserved by $U$, so $\langle\psi|\phi\rangle = \langle\psi|\phi\rangle^2$,
which forces $\langle\psi|\phi\rangle \in \{0, 1\}$.
\end{proof}

\bibliographystyle{unsrt}
\bibliography{refs}
\end{document}
